\section{Exact Divisor Local Factor And Boundary Remainder}

\textbf{Theorem 2 (Exact divisor local factor and boundary remainder).}
For every pair of squarefree prime wheels $W,Z$ with $2\mid W$, every integer
interval $[L,U)$ with $L\lt U$, and every $m\ge1$, over relaxed-weight
integers in that interval that are also divisible by the fixed integer $m$:

\begin{equation*}
\mathcal N_m[L,U)=\rho(m)\ell_m+E_m[L,U),
\qquad
|E_m[L,U)|\le R-1.
\end{equation*}

Before studying divisor averages, the comparison density must account for how
the tested divisor meets the two wheels. Let $W$ and $Z$ be squarefree with
$2\mid W$, and define

\begin{equation*}
\mathcal N_m[L,U)
=
\#\{n\in[L,U):m\mid n,\ \gcd(n,W)=1,\ \gcd(n+2,Z)=1\}.
\end{equation*}

Writing $n=mk$ converts the numerical interval to

\begin{equation*}
K_m[L,U)
=
\left[
\left\lceil\frac Lm\right\rceil,
\left\lceil\frac Um\right\rceil
\right)\cap\mathbb Z,
\qquad
\ell_m=|K_m[L,U)|.
\end{equation*}

If $\gcd(m,W)>1$, choose a prime $p\mid\gcd(m,W)$. Every $n=mk$ then has
$p\mid n$ and cannot be coprime to $W$. Thus

\begin{equation*}
\mathcal N_m[L,U)=0.
\end{equation*}

Adopt the convention $\rho(m)=E_m[L,U)=0$ whenever $\gcd(m,W)>1$: with this
convention the displayed identity
$\mathcal N_m[L,U)=\rho(m)\ell_m+E_m[L,U)$ holds trivially in this branch,
so the theorem's quantifier over every $m\ge1$ is exact rather than
implicitly restricted to the coprime case developed below.

Assume now $\gcd(m,W)=1$. For every prime $p\mid WZ$, let
$\lambda_p(m)$ count the allowed residues of $k$ modulo $p$. Direct local
analysis gives

\begin{equation*}
\lambda_p(m)=
\begin{cases}
p-1,&p\mid W,\ p\nmid Z,\\
1,&p=2,\ p\mid W,\ p\mid Z,\\
p-2,&p>2,\ p\mid W,\ p\mid Z,\\
p,&p\mid Z,\ p\nmid W,\ p\mid m,\\
p-1,&p\mid Z,\ p\nmid W,\ p\nmid m.
\end{cases}
\end{equation*}

Each row follows from an exact residue count:

\begin{itemize}
\item If $p\mid W$ and $p\nmid Z$, the hypothesis $\gcd(m,W)=1$ makes $m$
  invertible modulo $p$. The installed condition forbids only
  $k\equiv0\pmod p$, leaving $p-1$ classes.
\item If $p\mid W$ and $p\mid Z$, the two conditions forbid
  $k\equiv0\pmod p$ and $k\equiv-2m^{-1}\pmod p$. They coincide for $p=2$,
  leaving one class, and are distinct for $p>2$, leaving $p-2$ classes.
\item If $p\mid Z$ and $p\nmid W$, then $p$ is odd because $2\mid W$. When
  $p\mid m$, one has $mk+2\equiv2\not\equiv0\pmod p$ for every $k$, so all
  $p$ classes survive. When $p\nmid m$, invertibility of $m$ leaves exactly
  one forbidden class and therefore $p-1$ allowed classes.
\end{itemize}

These cases are exhaustive because every prime dividing $WZ$ belongs to the
installed wheel only, both wheels, or the relaxed wheel only. This proves the
local table.

Put

\begin{equation*}
R=\prod_{p\mid WZ}p,
\qquad
\rho(m)=\prod_{p\mid WZ}\frac{\lambda_p(m)}p.
\end{equation*}

CRT proves that every complete block of $R$ consecutive $k$ values contains
exactly $R\rho(m)$ allowed values.

Let the $k$ interval have length $\ell_m=qR+s$, with $0\le s\lt R$, and let
$C_m$ count allowed values in its final $s$ positions. Then

\begin{equation*}
\begin{aligned}
\mathcal N_m[L,U)
&=qR\rho(m)+C_m
&&[\text{Complete Periods Plus Remainder}]\\
&=\rho(m)\ell_m+
\left(C_m-s\rho(m)\right).
&&[\text{Substitution}]
\end{aligned}
\end{equation*}

Defining $E_m[L,U)=C_m-s\rho(m)$ gives

\begin{equation*}
\mathcal N_m[L,U)
=\rho(m)\ell_m+E_m[L,U),
\qquad
|E_m[L,U)|\le s\le R-1.
\qquad\blacksquare\ \text{[Q.E.D.]}
\end{equation*}

In the candidate range $Z\mid W$, all divisors coprime to $W$ have the same
density

\begin{equation*}
\rho_{Q,z}
=
\frac12
\prod_{2\lt p\lt z}\left(1-\frac2p\right)
\prod_{z\le p\lt Q}\left(1-\frac1p\right).
\end{equation*}

This is sieve dimension two below $z$ and dimension one from $z$ to $Q$.
The exact pointwise remainder bound is not a Type-I theorem because the
primorial $R$ is much larger than the square-safe interval.

No maintained Scala theorem currently models both squarefree wheels, all five
local cases, CRT composition, and the arbitrary interval remainder; we leave
it unverified rather than represent it with speculative code. Proving it
would require establishing the local table one prime at a time and then
using a verified CRT product lemma. The complete mathematical proof
is also recorded, with the same derivation, in \href{https://github.com/thiagomata/prime-numbers/blob/master/properties/sieve-sequence/relaxed-almost-prime-divisor-local-factor.md}{Relaxed Almost-Prime Weight Has An Exact Divisor Local
Factor} \cite{divisorlocalfactor2026}.
